% Generated from validated Article Draft Result. Do not hand-edit; revise the structured draft instead.
\section{Detailed Chronology――paper、release、accessを同じ日付にしない}
\label{pkg:annual-chronology}
\noindent\textbf{年次史では『その技術が2021年にあった』だけでは足りない。paperの初出、組織によるannouncement、model weightのrelease、APIやproductへのaccessは別のchronology stateである。本章はAI21 Studio / Jurassic-1、MT-NLG、GLIDE、Latent Diffusion Models、language-model risk taxonomyの資料を使い、本文で圧縮した2021年後半の動きを公開形態の違いごとに読み直す。}\autocite{src-20fd9bba2e7a,src-50244427230b,src-a4561c635fe6,src-b307b4b82867,src-bf581c89a17f}

Chronologyで最初に固定すべきなのは日付ではなくevent typeである。paperが公開されたこと、projectやvendorがfirst-party pageでmodelを告知したこと、利用者がAPIへaccessできること、product previewが始まったことは、相互に関連していても同一の出来事ではない。本号では、それらを一つの『release date』へ丸めず、research artifactとdistribution/access boundaryを分離して扱う。この原則は年次版ほど重要で、後から有名になった日付だけで一年を再構成することを防ぐ。\autocite{src-50244427230b,src-b307b4b82867,src-bf581c89a17f}


Source Intakeだけでは歴史上重要なfirst-party access boundaryが必ずしも同じ密度で残らないため、本号ではtargeted supplemental primary sourcesをchronologyの補助線に使っている。Jurassic-1 / AI21 Studioは独立providerによるAPI access、Megatron-Turing NLGはMicrosoft / NVIDIAによるfrontier-scale modelのfirst-party recordとして扱う。これらはpaper中心の研究chronologyとは異なる情報を補い、2021年にmodelの『存在』と『利用・公開surface』が分離していたことを示す。\autocite{src-b307b4b82867,src-bf581c89a17f}


年末側では、GLIDEとLatent Diffusion Modelsがdiffusion trajectoryを異なる方向へ延ばしている。前者はtext-guided generation/editing、後者はlatent-space efficiencyという別の設計点を持つため、単に『年末に画像生成が進歩した』とまとめるべきではない。paper artifactとしての初出と、後年それらが広いecosystemへ結び付いた時期も切り分ける。2021年のchronologyに置くのは、その技術的構成要素が当年の研究記録に現れたことまでである。\autocite{src-50244427230b,src-a4561c635fe6}


同じ年末の密度には、生成能力だけでなくevaluation/riskの成熟も含まれる。Ethical and social risks of harm from Language Modelsは、language-model riskを明示的なtaxonomyとして扱う研究記録であり、media generationやfrontier-scale modelと異なるtrajectoryに属する。年末に複数の重要artifactが集中して見えても、それを『2021年の本体は12月だった』という物語へ変換しない。むしろ一年を通じて分化したlayerが、後半に同時に可視化されたと読む。\autocite{src-20fd9bba2e7a,src-50244427230b,src-a4561c635fe6}


このDetailed Chronologyでは、Evidence Cardが正確なevent dateを固定していない項目に対して、現在の知識やURL表記から日付を補完しない。年次順序を理解する手掛かりと、厳密に確認済みのtimestampは別物だからである。paper publication、official announcement、open-weight release、API availability、product previewの区別を保ち、正確な日付が必要な箇所は一次資料で固定された場合にだけ記載する。この規律により、後年の検索結果から作られた擬似的な一日単位chronologyを避ける。\autocite{src-20fd9bba2e7a,src-b307b4b82867,src-bf581c89a17f}


\begin{claimboundary}
Chronologyに採用したfirst-party pageやpaperは、artifactの存在と本号で扱う公開境界を確認するための一次資料である。各組織・著者の性能評価を独立再現したものではなく、Evidenceで固定されていない正確な公開時刻を推定して補わない。\autocite{src-20fd9bba2e7a,src-50244427230b,src-a4561c635fe6,src-b307b4b82867,src-bf581c89a17f}
\end{claimboundary}
